\documentclass[10pt,a4paper]{amsart}
\usepackage[margin=19mm]{geometry}
\usepackage{amsmath,amssymb,mathtools}
\usepackage[scr=rsfs,cal=euler]{mathalfa}
\usepackage{xfrac}
\usepackage[unicode,hidelinks,bookmarks=false]{hyperref}
\newcommand{\ie}{i.e.\ }
\newcommand{\cf}{cf.\ }
\newcommand{\resp}{respectively\ }
\newcommand{\Subsection}{\subsection}
\providecommand{\nobreakdash}{\nobreak\hbox{-}}
\setlength{\emergencystretch}{2em}
\multlinegap0pt
\usepackage{amssymb}
\usepackage{tikz}
\usepackage[normalem]{ulem}
\newtheorem{proposition}{Proposition}
\def \N {{\mathbb N}}
\def\ge{\geqslant}
\def\geq{\geqslant}
\def\le{\leqslant}
\def\leq{\leqslant}
\title{Theorem $(1+1.9)$ on the Goldbach Conjecture}
\author{Jiamin Li, Jianya Liu}
\date{Source: arXiv:2606.05224v2 (CC BY 4.0)}
\begin{document}
\maketitle

\noindent\emph{Source and licence.} Theorem $(1+1.9)$ on the Goldbach Conjecture. Version arXiv:2606.05224v2 (CC BY 4.0), distributed by arXiv under the Creative Commons Attribution 4.0 International licence, http://creativecommons.org/licenses/by/4.0/, as recorded in the arXiv licence metadata for this exact version and shown on the abstract page https://arxiv.org/abs/2606.05224v2. Full licence notice: \texttt{licenses/arxiv-2606.05224-CC-BY-4.0.txt}.\par\smallskip

\noindent\textit{Editorial scope.} Counted unit: the article's proposition Goldbachweight2 (source0.tex lines 1374--1440). The article's complete original proof of it is delivered with it (source0.tex lines 1442--1645, one \texttt{proof} environment of 204 lines). No mathematics is written, repaired or re-derived: the statement and the proof are the article's own lines, in the article's own notation, and no retained line is substituted or re-flowed.  Retained uncounted as the source-local context the proof needs: lines 907--940; lines 958--961; lines 963--965; lines 979--981; lines 1254--1274.  Nothing was omitted from any retained span, so the package carries no omission record.  No retained line contains a citation, so the wrapper carries no bibliography and no bibliography entry.  No retained line contains a figure environment, an image-inclusion command or drawing code, so no image is delivered and the package declares no asset dependency.  Editorial formatting only, in addition: an AMS \texttt{amsart} preamble that restates the article's own declarations and macros used by the retained lines, namely the environments \texttt{proposition} on separate counters, together with the standard packages those lines need.  The retained environments print in this document as Proposition 1 in order of appearance, whereas the article prints the same items with the numbering of its own full document.  The complete native source of the article as distributed in the arXiv e-print for this exact version is bundled unchanged as \texttt{sources/202106/source0.tex}.
\begin{proposition}\label{Goldbachbig}
Let 
$$
1.5+3\varepsilon<a<2,\quad \tau=\dfrac{a-1}{a}-\varepsilon, \quad \frac{1}{21}<\kappa<\sigma\le \frac{1}{3}. 
$$ 
Then we have 
\begin{equation}\label{Goldbachbigw2}
\begin{split}
2D_{1,a}(N)
\ge\ & 2S_1(\kappa)-2S_2(\tau)-S_3(\kappa,\sigma)-2S_4(\sigma)\\
&-S_5(\kappa,\sigma)+S_6(\kappa,\sigma)+O(N^{1-\kappa}),
\end{split}
\end{equation}
where
\[
S_1(\kappa) := S(\mathscr{A}; \mathscr{P}(N), N^\kappa), 
\]
\[
S_2(\kappa) :=\sum\limits_{\substack{N^\kappa\le p\le N^{1/2} \\(p, N)=1}}S(\mathscr{A}_p ; \mathscr{P}(N), p), 
\]
\[
S_3(\kappa, \sigma) := \sum_{\substack{N^\kappa \le p \le N^\sigma \\ (p, N) = 1}} S(\mathscr{A}_p; \mathscr{P}(N), N^\kappa),
\] 
\[
S_4(\sigma) := \sum_{\substack{N^\sigma \le p_1 \le p_2 \le (N/p_1)^{1/2} \\ (p_1 p_2, N) = 1}} S(\mathscr{A}_{p_1 p_2}; \mathscr{P}(N p_1), p_2), 
\]
\[
S_5(\kappa, \sigma) := \sum_{\substack{N^\kappa \le p_1 \le N^\sigma \le p_2 \le (N/p_1)^{1/2} \\ (p_1 p_2, N) = 1}}  S(\mathscr{A}_{p_1 p_2}; \mathscr{P}(N p_1), p_2),
\]
and 
\[
S_6(\kappa, \sigma) := \sum_{\substack{N^\kappa \le p_1 \le p_2 \le p_3 \le N^\sigma \\ (p_1 p_2 p_3, N) = 1}} S(\mathscr{A}_{p_1 p_2 p_3};  \mathscr{P}(N p_1), p_2).
\]
\end{proposition}

\begin{equation}\label{Buchstab1}
\sum_{\substack{N^\kappa \le p \le N^{1/2} \\ (p, N) = 1}} S(\mathscr{A}_{p^2} ; \mathscr{P}(N), p) 
\ll N^{1-\kappa} 
\end{equation}

\begin{equation}\label{Buchstab0} 
S(\mathscr{A}_{p} ; \mathscr{P}(N), p) \ll N^{1-\kappa}
\end{equation}

		\begin{equation}\label{Buchstab2}
		\sum_{\substack{N^\sigma \le p \le N^{1/2} \\ (p, N) = 1}} S\left(\mathscr{A}_{p^2} ; \mathscr{P}(N), p\right) \ll N^{1-\sigma}.
		\end{equation}

\begin{proposition}\label{Goldbachsmall}
Let 
$$
1.05+3\varepsilon<a<1.5,\quad \tau=\dfrac{a-1}{a}-\varepsilon, \quad 
\frac{1}{21}<\kappa<\sigma\le \tau< \frac{1}{3}. 
$$ 
Then we have
		\begin{equation}\label{Goldbachsmallw2}
		\begin{split}
		2D_{1,a}(N)\ge\ & 2S_1(\kappa)-2S_2(\tau)-S_3(\kappa,\sigma)-S_5(\kappa,\sigma)+S_6(\kappa,\sigma)\\
		&-2S_7(\sigma,\tau)-2S_8(\sigma,\tau)+O(N^{1-\kappa}),
		\end{split}
		\end{equation}
where $S_1, \ldots, S_6$ are as in Proposition~\ref{Goldbachbig}, and 
\begin{equation}
\begin{split}
S_{7}(\sigma,\tau):=&\sum_{\substack{N^\sigma \le p_1 \le p_2 \le N^\tau \\ (p_1 p_2, N) = 1}}  S(\mathscr{A}_{p_1 p_2}; \mathscr{P}(N p_1), p_2), \\
S_{8}(\sigma,\tau):=&\sum_{\substack{N^\sigma \le p_1 \le N^\tau\le p_2 \le (N/p_1)^{1/2} \\ (p_1 p_2, N) = 1}}  S(\mathscr{A}_{p_1 p_2}; \mathscr{P}(N p_1), p_2). 
\end{split}
\end{equation}
\end{proposition}

\begin{proposition}\label{Goldbachweight2}
Let 
$$
1.05+3\varepsilon<a<1.5,\quad \tau=\dfrac{a-1}{a}-\varepsilon, \quad \dfrac{1}{18}<\alpha<\beta<\gamma<\delta< \tau<\dfrac{1}{3} 
$$ 
such that
\begin{equation}\label{Con/bcd}
\beta+\gamma+2\delta< 1. 
\end{equation}
Then 
\begin{equation}\label{Goldbachsmallw4}
		\begin{split}
		4D_{1,a}(N) \geq\ &3H_1 + H_2 -4H_3- H_4- H_5 + H_6 + H_7- 2H_8 \\
		& - 2H_9 -2H_{10}- 2H_{11} - H_{12} - H_{13}-H_{14}-H_{15} \\ 
		& -H_{16}-H_{17} + O(N^{1 - \alpha}),
		\end{split}
		\end{equation}
		where
\[
H_1 := S(\mathscr{A}; \mathscr{P}(N), N^{\alpha}),\quad H_2 := S(\mathscr{A}; \mathscr{P}(N), N^{\beta}),
\]
\[
H_3 :=\sum\limits_{\substack{N^\tau\le p\le N^{1/2} \\(p, N)=1}}S(\mathscr{A}_p ; \mathscr{P}(N), p),
\quad 
H_4 := \sum_{\substack{N^{\alpha} \leq p \le N^{\delta} \\ (p, N) = 1}} S(\mathscr{A}_p; \mathscr{P}(N), N^{\alpha}), 
\]
\[
H_5 := \sum_{\substack{N^{\alpha} \leq p \le N^{\gamma} \\ (p, N) = 1}} S(\mathscr{A}_p; \mathscr{P}(N), N^{\alpha}),
\]
\[ 
H_6 := \sum_{\substack{N^{\alpha} \leq p_1 \le p_2 \le N^{\beta} \\ (p_1 p_2, N) = 1}} S(\mathscr{A}_{p_1 p_2}; \mathscr{P}(N), N^{\alpha}),
\]
\[ 
H_7 := \sum_{\substack{N^{\alpha} \leq p_1 \le N^{\beta} \le p_2 \le N^{\gamma} \\ (p_1 p_2, N) = 1}} S(\mathscr{A}_{p_1 p_2}; \mathscr{P}(N), N^{\alpha}),
\] 
\[
H_8 := \sum_{\substack{N^{\gamma} \leq p_1 \le p_2 \le N^{\tau} \\ (p_1 p_2, N) = 1}} S(\mathscr{A}_{p_1 p_2}; \mathscr{P}(N p_1), p_2),
\]
\[ 
H_{9} := \sum_{\substack{N^{\delta} \leq p_1 \le p_2 \le N^{\tau} \\ (p_1 p_2, N) = 1}} S(\mathscr{A}_{p_1 p_2}; \mathscr{P}(N p_1), p_2),
\]
\[
H_{10} := \sum_{\substack{N^{\gamma} \leq p_1 \le N^{\tau}\le p_2 \le (N/p_1)^{1/2} \\ (p_1 p_2, N) = 1}} S(\mathscr{A}_{p_1 p_2}; \mathscr{P}(N p_1), p_2),
\]
\[
H_{11} := \sum_{\substack{N^{\delta} \leq p_1 \le N^{\tau}\le p_2 \le (N/p_1)^{1/2} \\ (p_1 p_2, N) = 1}} S(\mathscr{A}_{p_1 p_2}; \mathscr{P}(N p_1), p_2),
\] 
\[
H_{12} := \sum_{\substack{N^{\alpha} \leq p_1 \le N^{\delta}\leq p_2 \le (N/p_1)^{1/2} \\ (p_1 p_2, N) = 1}} S(\mathscr{A}_{p_1 p_2}; \mathscr{P}(N p_1), p_2),
\] 
\[
H_{13} := \sum_{\substack{N^{\beta} \leq p_1 \le N^{\gamma} \leq p_2 \le (N/p_1)^{1/2} \\ (p_1 p_2, N) = 1}} S(\mathscr{A}_{p_1 p_2}; \mathscr{P}(N p_1), (N/p_1 p_2)^{1/2}),
\]
\[
H_{14} := \sum_{\substack{N^\beta \le p_1\le N^{\gamma}\\N^\delta \le p_2\le p_3\le (N/p_1p_2)^{1/2} \\ (p_1 p_2 p_3, N) = 1}}  S(\mathscr{A}_{p_1 p_2 p_3}; \mathscr{P}(Np_1), p_2),
\]
\[ 
H_{15} := \sum_{\substack{N^\beta \le p_1\le N^\gamma \le p_2\le N^{\delta}\le p_3\le (N/p_1p_2)^{1/2} \\ (p_1 p_2 p_3, N) = 1}}  S(\mathscr{A}_{p_1 p_2 p_3}; \mathscr{P}(Np_1), p_2),
\]
\[
H_{16} := \sum_{\substack{N^{\alpha} \leq p_1 \le p_2 \le p_3 \le p_4 \le N^{\beta} \\ (p_1 p_2 p_3 p_4, N) = 1}} S(\mathscr{A}_{p_1 p_2 p_3 p_4}; \mathscr{P}(Np_1), p_2),
\]
and 
\[
H_{17} := \sum_{\substack{N^{\alpha} \leq p_1 \le p_2 \le p_3 \le N^{\beta} \leq p_4 \le N^{\gamma}\\ (p_1 p_2 p_3 p_4, N) = 1}} S(\mathscr{A}_{p_1 p_2 p_3 p_4}; \mathscr{P}(Np_1), p_2).
\] 
\end{proposition}

\begin{proof}
First, we apply Proposition \ref{Goldbachsmall} with $(\kappa,\sigma,a)=(\alpha,\delta,a)$ to obtain
\begin{equation}\label{GoldW1}
\begin{split}
2D_{1,a}(N)\ge\ & 2S_1(\alpha)-2S_2(\tau)-S_3(\alpha,\delta)-S_5(\alpha,\delta)+S_6(\alpha,\delta)\\
&-2S_7(\delta,\tau)-2S_8(\delta,\tau)+O(N^{1-\alpha})\\
=\ & 2H_1-2H_3-H_4-H_{12}+S_6(\alpha,\delta)\\ 
& -2H_9-2H_{11}+O(N^{1-\alpha}). 
\end{split}
\end{equation}
Another application of Proposition \ref{Goldbachsmall} with $(\kappa,\sigma,a)=(\beta,\gamma,a)$ gives 
\begin{equation}\label{GoldW2}
\begin{split}
2D_{1,a}(N)
\ge\ & 2S_1(\beta)-2S_2(\tau)-S_3(\beta,\gamma)-S_5(\beta,\gamma)+S_6(\beta,\gamma)\\
& -2S_7(\gamma,\tau)-2S_8(\gamma,\tau)+O(N^{1-\beta})\\
=\ & 2H_2-2H_3-S_3(\beta,\gamma)-S_5(\beta,\gamma)+S_6(\beta,\gamma) \\ 
& -2H_8-2H_{10}+O(N^{1-\beta}).
\end{split}
\end{equation}
Adding \eqref{GoldW1} and \eqref{GoldW2}, we get 
\begin{equation}\label{GoldW3}
\begin{split}
4D_{1,a}(N)\ge\ & 2H_1+2H_2-4H_3-H_4-2H_8-2H_9-2H_{10}-2H_{11}\\
& -H_{12} -S_3(\beta,\gamma)-S_5(\beta,\gamma)+S_6(\alpha,\delta)+O(N^{1-\alpha}),
\end{split}
\end{equation}
where we have used the assumption $\alpha<\beta$, and have dropped the term $S_6(\beta,\gamma)$ by non-negativity. 

Next we are going to analyze the terms on the right-hand side. In the following Buchstab's identity as well as 
estimates like \eqref{Buchstab1}, \eqref{Buchstab0} and \eqref{Buchstab2} will be used repeatedly, and this will not be 
pointed out at each circumstances. 
The $H_2$ on the right-hand side of \eqref{GoldW3} can be written as
\begin{equation*}
\begin{split}
H_2=\ &S(\mathscr{A};\mathscr{P}(N),N^\beta)\\
=\ & S(\mathscr{A};\mathscr{P}(N),N^\alpha)-\sum\limits_{\substack{N^\alpha\le p\le N^{\beta} \\(p, N)=1}}S(\mathscr{A}_p ; \mathscr{P}(N), p)\\
=\ & S(\mathscr{A};\mathscr{P}(N),N^\alpha)
-\sum\limits_{\substack{N^\alpha\le p\le N^{\beta} \\(p, N)=1}}S(\mathscr{A}_p; \mathscr{P}(N), N^\alpha) \\ 
&+\sum_{\substack{N^\alpha \le p_1 \le p_2 \le N^\beta \\ (p_1 p_2, N) = 1}}  S(\mathscr{A}_{p_1 p_2}; \mathscr{P}(N), p_1)+O(N^{1-\alpha}), 
\end{split}
\end{equation*}
and this can be further transform as 
\begin{equation}\label{GoldW4}
\begin{split}
H_2=\ & S(\mathscr{A};\mathscr{P}(N),N^\alpha)-\sum\limits_{\substack{N^\alpha\le p\le N^{\beta} \\(p, N)=1}}S(\mathscr{A}_p ; \mathscr{P}(N), N^\alpha)\\
		&+\sum_{\substack{N^\alpha \le p_1 \le p_2 \le N^\beta \\ (p_1 p_2, N) = 1}}  S(\mathscr{A}_{p_1 p_2}; \mathscr{P}(N), N^\alpha)\\
		&-\sum_{\substack{N^\alpha \le p_1 \le p_2 \le p_3\le N^\beta \\ (p_1 p_2 p_3, N) = 1}}  S(\mathscr{A}_{p_1 p_2 p_3}; \mathscr{P}(N), p_1)+O(N^{1-\alpha})\\
=\ &H_1-H_{18}+H_6-H_{19}+O(N^{1-\alpha}),
\end{split}
\end{equation}
where 
\begin{equation}
\begin{split}
H_{18}:=\ &\sum\limits_{\substack{N^\alpha\le p\le N^{\beta} \\(p, N)=1}}S(\mathscr{A}_p ; \mathscr{P}(N), N^\alpha),\\
H_{19}:=\ &\sum_{\substack{N^\alpha \le p_1 \le p_2 \le p_3\le N^\beta \\ (p_1 p_2 p_3, N) = 1}}  S(\mathscr{A}_{p_1 p_2 p_3}; \mathscr{P}(N), p_1). 
\end{split}
\end{equation}
We replace one of the $H_2$'s on the right-hand side of \eqref{GoldW3} by \eqref{GoldW4}, so that 
 \begin{equation}\label{GoldW5}
\begin{split}
4D_{1,a}(N)\ge\ & 3H_1+H_2-4H_3-H_4+H_6-2H_8\\
& -2H_9-2H_{10}-2H_{11} -H_{12} -H_{18}-H_{19} \\ 
& -S_3(\beta,\gamma)-S_5(\beta,\gamma)+S_6(\alpha,\delta)+O(N^{1-\alpha}). 
\end{split}
\end{equation}

For the term $S_3(\beta,\gamma)$ on the right-hand side, 
we have that 
\begin{equation*}
\begin{split}
S_3(\beta,\gamma)=\ &\sum_{\substack{N^\beta \le p \le N^\gamma \\ (p, N) = 1}} S(\mathscr{A}_p; \mathscr{P}(N), N^\beta)\\
=\ & \sum_{\substack{N^\beta \le p \le N^\gamma \\ (p, N) = 1}} S(\mathscr{A}_p; \mathscr{P}(N), N^\alpha)-\sum_{\substack{N^{\alpha} \leq p_1 \le N^{\beta} \le p_2 \le N^{\gamma} \\ (p_1 p_2, N) = 1}} S(\mathscr{A}_{p_1 p_2}; \mathscr{P}(N), p_1), 
\end{split}
\end{equation*}
and hence 
\begin{equation}\label{GoldW6}
\begin{split}
S_3(\beta,\gamma)
=\ & \sum_{\substack{N^\alpha \le p \le N^\gamma \\ (p, N) = 1}} S(\mathscr{A}_p; \mathscr{P}(N), N^\alpha)-\sum_{\substack{N^\alpha \le p \le N^\beta \\ (p, N) = 1}} S(\mathscr{A}_p; \mathscr{P}(N), N^\alpha)\\
		&-\sum_{\substack{N^{\alpha} \leq p_1 \le N^{\beta} \le p_2 \le N^{\gamma} \\ (p_1 p_2, N) = 1}} S(\mathscr{A}_{p_1 p_2}; \mathscr{P}(N), N^{\alpha})\\
		&+\sum_{\substack{N^{\alpha} \leq p_1 \le p_2 \le N^{\beta} \leq p_3 \le N^{\gamma}\\ (p_1 p_2 p_3, N) = 1}} S(\mathscr{A}_{p_1 p_2 p_3}; \mathscr{P}(N), p_1)+O(N^{1-\alpha})\\
		=\ &H_5-H_{18}-H_{7}+H_{20}+O(N^{1-\alpha}),
\end{split}
\end{equation}
where
\begin{equation}
H_{20}:=\sum_{\substack{N^{\alpha} \leq p_1 \le p_2 \le N^{\beta} \leq p_3 \le N^{\gamma}\\ (p_1 p_2 p_3, N) = 1}} 
S(\mathscr{A}_{p_1 p_2 p_3}; \mathscr{P}(N), p_1). 
\end{equation}
Inserting \eqref{GoldW6} into \eqref{GoldW5} yields 
\begin{equation}\label{GoldW7}
\begin{split}
4D_{1,a}(N)
\ge\ & 3H_1+H_2-4H_3-H_4-H_5+H_6+H_7 -2H_8\\ 
& -2H_9-2H_{10}-2H_{11} -H_{12}-H_{19}-H_{20} \\ 
& -S_5(\beta,\gamma)+S_6(\alpha,\delta)+O(N^{1-\alpha}).
\end{split}
\end{equation}

For the term $S_5(\beta,\gamma)$ on the right-hand side, 
we have similarly that 
		\begin{equation}\label{GoldW8}
		\begin{split}
		S_5(\beta,\gamma)=\ &\sum_{\substack{N^\beta \le p_1 \le N^\gamma \le p_2 \le (N/p_1)^{1/2} \\ (p_1 p_2, N) = 1}}  S(\mathscr{A}_{p_1 p_2}; \mathscr{P}(N p_1), p_2)\\
		=\ & \sum_{\substack{N^{\beta} \leq p_1 \le N^{\gamma} \leq p_2 \le (N/p_1)^{1/2} \\ (p_1 p_2, N) = 1}} S(\mathscr{A}_{p_1 p_2}; \mathscr{P}(N p_1), (N/p_1 p_2)^{1/2})\\
		&+\sum_{\substack{N^{\beta} \leq p_1 \le N^{\gamma} \leq p_2 \le p_3\le (N/p_1p_2)^{1/2} \\ (p_1 p_2 p_3, N) = 1}} S(\mathscr{A}_{p_1 p_2}; \mathscr{P}(N p_1), p_3)+O(N^{1-\beta})\\
		=\ &H_{13}+H_{21}+O(N^{1-\alpha}),
		\end{split}
		\end{equation}
		where 
		\begin{equation}
		H_{21}:=\sum_{\substack{N^{\beta} \leq p_1 \le N^{\gamma} \leq p_2 \le p_3\le (N/p_1p_2)^{1/2} \\ (p_1 p_2 p_3, N) = 1}} S(\mathscr{A}_{p_1 p_2}; \mathscr{P}(N p_1), p_3). 
		\end{equation}
Inserting \eqref{GoldW8} into \eqref{GoldW7}, we have
\begin{equation}\label{GoldW9}
\begin{split}
4D_{1,a}(N)
\ge\ & 3H_1+H_2-4H_3-H_4-H_5+H_6+H_7-2H_8\\
&-2H_9-2H_{10}-2H_{11}-H_{12}-H_{13}-H_{19}-H_{20}-H_{21} \\ 
& +S_6(\alpha,\delta)+O(N^{1-\alpha}).
\end{split}
\end{equation}
		
For the term $S_6(\alpha, \delta)$ on the right-hand side, 
we have similarly that 
		\begin{equation}\label{GoldW10}
		\begin{split}
		S_6(\alpha, \delta) =& \sum_{\substack{N^\alpha \le p_1 \le p_2 \le p_3 \le N^{\delta} \\ (p_1 p_2 p_3, N) = 1}} S(\mathscr{A}_{p_1 p_2 p_3}; \mathscr{P}(N p_1), p_2)\\
		=&\sum_{\substack{N^\alpha \le p_1 \le p_2 \le p_3\le N^\beta \\ (p_1 p_2 p_3, N) = 1}}  S(\mathscr{A}_{p_1 p_2 p_3}; \mathscr{P}(Np_1), p_2)\\
		&+\sum_{\substack{N^\alpha \le p_1 \le p_2\le N^\beta \le p_3\le N^{\delta} \\ (p_1 p_2 p_3, N) = 1}}  S(\mathscr{A}_{p_1 p_2 p_3}; \mathscr{P}(Np_1), p_2)\\
		&+\sum_{\substack{N^\alpha \le p_1 \le N^\beta\le p_2\le p_3\le N^{\delta} \\ (p_1 p_2 p_3, N) = 1}}  S(\mathscr{A}_{p_1 p_2 p_3}; \mathscr{P}(Np_1), p_2)\\
		&+\sum_{\substack{N^\beta \le p_1 \le p_2\le p_3\le N^{\delta} \\ (p_1 p_2 p_3, N) = 1}}  S(\mathscr{A}_{p_1 p_2 p_3}; \mathscr{P}(Np_1), p_2). 
\end{split}
\end{equation}
The next-to-last term is dropped by non-negativity. The last term above is, by \eqref{Con/bcd}, 
\begin{equation}\label{GoldW11}
		\begin{split}
%		&\sum_{\substack{N^\beta \le p_1 \le p_2\le p_3\le N^{\delta} \\ (p_1 p_2 p_3, N) = 1}}  S(\mathscr{A}_{p_1 p_2 p_3}; \mathscr{P}(Np_1), p_2)\\
		& \ge \sum_{\substack{N^\beta \le p_1\le N^\gamma \le p_2\le p_3\le N^{\delta} \\ (p_1 p_2 p_3, N) = 1}}  S(\mathscr{A}_{p_1 p_2 p_3}; \mathscr{P}(Np_1), p_2)\\
		& = \sum_{\substack{N^\beta \le p_1\le N^\gamma \le p_2\le p_3\le (N/p_1p_2)^{1/2} \\ (p_1 p_2 p_3, N) = 1}}  S(\mathscr{A}_{p_1 p_2 p_3}; \mathscr{P}(Np_1), p_2)\\
		 &\quad  -\sum_{\substack{N^\beta \le p_1\le N^{\gamma}\\N^\delta \le p_2\le p_3\le (N/p_1p_2)^{1/2} \\ (p_1 p_2 p_3, N) = 1}}  S(\mathscr{A}_{p_1 p_2 p_3}; \mathscr{P}(Np_1), p_2)\\
		&\quad - \sum_{\substack{N^\beta \le p_1\le N^\gamma \le p_2\le N^{\delta}\le p_3\le (N/p_1p_2)^{1/2} \\ (p_1 p_2 p_3, N) = 1}}  S(\mathscr{A}_{p_1 p_2 p_3}; \mathscr{P}(Np_1), p_2)\\
		&= H_{22}-H_{14}-H_{15}, 
\end{split}
\end{equation}
where
\begin{equation}
H_{22}:=\sum_{\substack{N^{\beta} \leq p_1 \le N^{\gamma} \leq p_2 \le p_3\le (N/p_1p_2)^{1/2} \\ (p_1 p_2 p_3, N) = 1}} S(\mathscr{A}_{p_1 p_2}; \mathscr{P}(N p_1), p_2). 
\end{equation}
It follows from from \eqref{GoldW10} and \eqref{GoldW11} that 
\begin{equation}\label{GoldW12}
			S_6(\alpha,\delta)\ge H_{23}+H_{24}+H_{22}-H_{14}-H_{15},
		\end{equation}
		where
		\begin{equation}
		\begin{split}
		H_{23}:=&\sum_{\substack{N^\alpha \le p_1 \le p_2 \le p_3\le N^\beta \\ (p_1 p_2 p_3, N) = 1}}  S(\mathscr{A}_{p_1 p_2 p_3}; \mathscr{P}(Np_1), p_2),\\
		H_{24}:=&\sum_{\substack{N^\alpha \le p_1 \le p_2\le N^\beta \le p_3\le N^{\delta} \\ (p_1 p_2 p_3, N) = 1}}  S(\mathscr{A}_{p_1 p_2 p_3}; \mathscr{P}(Np_1), p_2).
		\end{split}
		\end{equation}
Inserting \eqref{GoldW12} into \eqref{GoldW9}, one gets 
		\begin{equation}\label{GoldW13}
		\begin{split}
		4D_{1,a}(N)
		\ge\ 
		& 3H_1+H_2-4H_3-H_4-H_5+H_6+H_7-2H_8\\
		&-2H_9-2H_{10}-2H_{11}-H_{12}-H_{13}-H_{14}-H_{15}\\
		&-H_{19}-H_{20}-H_{21}+H_{22}+H_{23}+H_{24}+O(N^{1-\alpha}).
		\end{split}
		\end{equation}
		
Now we only need to prove  
\begin{equation}\label{GoldW14}
H_{22}+H_{23}+H_{24}\ge H_{19}+H_{20}+H_{21}-H_{16}-H_{17},
\end{equation}
so that \eqref{Goldbachsmallw4} holds by \eqref{GoldW13} and \eqref{GoldW14}.
		
By Buchstab's identity, we have
\begin{equation}\label{GoldW15}
\begin{split}
H_{19}-H_{16}
=\ &\sum_{\substack{N^\alpha \le p_1 \le p_2 \le p_3\le N^\beta \\ (p_1 p_2 p_3, N) = 1}}  S(\mathscr{A}_{p_1 p_2 p_3}; \mathscr{P}(N), p_1)\\&-\sum_{\substack{N^\alpha \le p_1 \le p_2 \le p_3\le p_4\le N^\beta \\ (p_1 p_2 p_3 p_4, N) = 1}}  S(\mathscr{A}_{p_1 p_2 p_3}; \mathscr{P}(Np_1), p_2)\\
=\ &\sum_{\substack{N^\alpha \le p_1 \le p_2 \le p_3\le N^\beta \\ (p_1 p_2 p_3, N) = 1}}  S(\mathscr{A}_{p_1 p_2 p_3}; \mathscr{P}(Np_1), p_2)+O(N^{1-\alpha})\\
=\ &H_{23}+O(N^{1-\alpha}). 
\end{split}
\end{equation}
Similarly, 
\begin{equation}\label{GoldW16}
\begin{split}
H_{20}-H_{17}
=\ &\sum_{\substack{N^{\alpha} \leq p_1 \le p_2 \le N^{\beta} \leq p_3 \le N^{\gamma}\\ (p_1 p_2 p_3, N) = 1}} S(\mathscr{A}_{p_1 p_2 p_3}; \mathscr{P}(N), p_1)\\&-\sum_{\substack{N^{\alpha} \leq p_1 \le p_2 \le p_3 \le N^{\beta} \leq p_4 \le N^{\gamma} \\ (p_1 p_2 p_3 p_4, N) = 1}} S(\mathscr{A}_{p_1 p_2 p_3 p_4}; \mathscr{P}(Np_1), p_2)\\
=\ &\sum_{\substack{N^{\alpha} \leq p_1 \le p_2 \le N^{\beta} \leq p_3 \le N^{\gamma}\\ (p_1 p_2 p_3, N) = 1}} S(\mathscr{A}_{p_1 p_2 p_3}; \mathscr{P}(Np_1), p_2)+O(N^{1-\alpha})\\
\le \ &H_{24}+O(N^{1-\alpha}). 
\end{split}
\end{equation}
Trivially, 
\begin{equation}\label{GoldW17}
\begin{split}
H_{21}=\sum_{\substack{N^{\beta} \leq p_1 \le N^{\gamma} \leq p_2 \le p_3\le (N/p_1p_2)^{1/2} \\ (p_1 p_2 p_3, N) = 1}} S(\mathscr{A}_{p_1 p_2}; \mathscr{P}(N p_1), p_3)\le H_{22}. 
\end{split}
\end{equation}
Now \eqref{GoldW14} follows from \eqref{GoldW15}, \eqref{GoldW16} and \eqref{GoldW17}.
\end{proof}

\end{document}
